\section{Analyse des résultats, complement}
\label{sec:analyse}

\subsection{DMA}

Bien que nous aillons traite les donnees en Sec. \ref{sec:DMAresultats}, le traitement permet de determiner plus precisement le pic de chacune des courbes en Fig. \ref{fig:DMA_resultats}. Couple a une autre fonction logiciel (recherche des extremeums locaux lors de changement de variation d'une fonction continue). Nous obtenons en compilant tous les extremumsle tableau suivant :

\\
\begin{table}[!ht]
\centering
\footnotesize
\setlength{\tabcolsep}{4pt}
\renewcommand{\arraystretch}{0.9}
\begin{tabular}{|c|c|c|}
\hline
	\textbf{Frequence [Hz]} 
	& \textbf{Peak [$^\circ$]}
	& \textbf{incertitude [$^\circ$]}
\\\hline
	1.000 & 140.5 & 1.606 \\
	2.000 & 143.5 & 1.594 \\
	5.000 & 148.6 & 1.601 \\
	10.000 & 150.5 & 1.616 \\
	20.000 & 155.7 & 1.518 \\
\hline
\end{tabular}
	\caption{Valeurs de temperatures auxquelles E (module d'Young) est maximal}
\label{tab:peaks}
\end{table}


On peut donc extraire tous les extremums locaux de Fig. \ref{fig:DMA_resultats} que l'on peut retrouver dans le tableau \ref{tab:peaks}. Grace a la relation Eq. \ref{eq:arrhenius} qu'on linearise tel que Eq. \ref{eq:arrheniuslin}, nous obtenons le graph d'Arrenhius suivant : 

\begin{figure}[hct]
        \centering
        \includegraphics[scale=0.40]{Seite_arrhenius_graph.pdf}
        \label{fig:arrenius}
        \caption{Arrhenius graph}
\end{figure}

Celui nous permet, comme nous le mentionnons en Sec. \ref{sec:theorie}, d'extraire l'energie d'activation $E_a$ associee a la transition vitreuse de notre echantillon de PMMA, soit\[
	\hspace{8cm} $E_a$ = 294.381 [kJ/mol] (87.6-159.8$^\circ$).
\]


